\originalchapter{研究生量子理论二：辐射、表示与散射}
\label{ch:graduate322}
Washington Taylor的8.322（Spring 2003）以8.321和8.07为先修，每周两次90分钟讲课和一次60分钟习题课，成绩完全来自作业。实际公开四个扫描讲义合本，共168页：含时微扰与辐射51页、对称与全同粒子53页、散射30页、相对论量子力学34页。另有3页课历和12份作业；公开导航没有考试或官方解答。这一学期不能借用另一学期试卷补成所谓官方考核。

本章按四个合本的主题连续重编，原已在本科篇详细展开的微扰、角动量及基础散射以章名回指，补齐多极跃迁、寿命、群表示、散射极点和Dirac结构。扫描的识别文本仅辅助定位，公式按物理定义重新推导；不将OCR乱码视为数学原稿。讲义中的Sakurai等教材引用保留信用，不附商书；全部新增答案为自主解答。

\originalsection{从相干跃迁到黄金规则}
相互作用绘景的Dyson级数见上一分编。对常值弱微扰从0作用至$T$，一阶跃迁概率为
\[
 P_{i\to f}(T)=\frac{4|V_{fi}|^2}{\hbar^2}
 \frac{\sin^2(\omega_{fi}T/2)}{\omega_{fi}^2}.
\]
对孤立二能级，这是有限时间的相干振荡近似，不能把它除以$T$后当成永久衰变率。若末态形成连续谱且耦合、态密度在宽度$\hbar/T$内平滑，分布极限$4\sin^2(\omega T/2)/(T\omega^2)\to2\pi\delta(\omega)$给
$\Gamma=(2\pi/\hbar)|V_{fi}|^2\rho(E_i)$。
需要同时有足够长的谱分辨时间和足够短的耗尽时间。原扫描第16--19与20--23页部分重复；同一推导只编一次，重复不计新增理论。

周期场$V(t)=We^{-i\omega t}+W^\dagger e^{i\omega t}$分别选$E_f-E_i=\pm\hbar\omega$。半经典场给吸收和受激发射，却没有真空中的自发发射。磁共振的精确旋转系解已在上一章给出；氨分子的两个宇称相反低能态也形成二能级系统，电偶极是奇宇称，故对角矩元为零、异宇称矩元可非零。微波场由这一非对角矩元驱动，不必把分子每个内部自由度都纳入两能级模型。

\originalsection{场量子化、偶极辐射与多极展开}
SI下取Coulomb规范$\nabla\cdot\vct A=0$。每个横向模式是频率$\omega_k=c|\vct k|$的谐振子，体积$V$中的规范化场为
\[
 \vct A=\sum_{\vct k\lambda}\sqrt{\frac{\hbar}{2\epsilon_0\omega_kV}}
 \vct\epsilon_{\vct k\lambda}
 (a_{\vct k\lambda}e^{i\vct k\cdot\vct r}
 +a^\dagger_{\vct k\lambda}e^{-i\vct k\cdot\vct r}),\quad
 H_\gamma=\sum\hbar\omega_k(a^\dagger a+1/2).
\]
模式满足$\vct k\cdot\vct\epsilon=0$，两个独立偏振。$a$与$a^\dagger$矩元的平方分别是$N$和$N+1$，正是吸收、受激发射、自发发射的区别。零点总能量形式发散，此处只计算有限跃迁矩元，未声称解决真空能的重整化。

长波$kr_{\rm atom}\ll1$下，最小耦合$-(q/m)\vct A\cdot\vct p$可用$\vct p=(im/\hbar)[H_0,\vct r]$转换成电偶极$-\vct d\cdot\vct E$。态密度与偏振求和给
\[
 \frac{\dd\Gamma_{E1}}{\dd\Omega}
 =\frac{\omega^3}{8\pi^2\epsilon_0\hbar c^3}
 (|\vct d_{fi}|^2-|\hat n\cdot\vct d_{fi}|^2),\qquad
 \Gamma_{E1}=\frac{\omega^3|\vct d_{fi}|^2}{3\pi\epsilon_0\hbar c^3}.
\]
角积分用$\int n_in_j\dd\Omega=4\pi\delta_{ij}/3$。线性$z$偶极分布为$\sin^2\theta$，圆偶极为$(1+\cos^2\theta)/2$，两者积分权重相同。电偶极秩1、奇宇称给$|j_i-j_f|\le1$且$0\not\to0$，宇称相反；对无自旋中心势为$\Delta l=\pm1$。具体氢$2p\to1s$径向积分及精细/超精细分支在后续PS3中计算。

若E1禁止，必须实际分解长波耦合，不能只列M1/E2结果。先取$H_0=p^2/(2m)+V(\vct r)$、横向单模$\vct A=A_0\vct\epsilon e^{i\vct k\cdot\vct r-i\omega t}$、$\vct\epsilon\cdot\vct k=0$；忽略二光子项$A^2$，线性耦合为$-(q/m)\vct A\cdot p$及自旋磁矩项。展开空间指数到一阶，其中
\[
 (\vct k\cdot r)(\vct\epsilon\cdot p)
 =\frac12[(\vct k\cdot r)(\vct\epsilon\cdot p)-(\vct\epsilon\cdot r)(\vct k\cdot p)]
 +\frac12[(\vct k\cdot r)(\vct\epsilon\cdot p)+(\vct\epsilon\cdot r)(\vct k\cdot p)].
\]
第一半是$(\vct k\times\vct\epsilon)\cdot L/2$，从而轨道磁偶极为$qL/(2m)$。第二半用位置二次算符的对易式
\[
 [H_0,(\vct k\cdot r)(\vct\epsilon\cdot r)]
 =-\frac{i\hbar}m[(\vct k\cdot p)(\vct\epsilon\cdot r)+(\vct k\cdot r)(\vct\epsilon\cdot p)]
\]
化成初末能量差的矩元；按发射/吸收的对应场相位取$E_f-E_i=\mp\hbar\omega$，得到长度形式$-q r_ir_j\partial_jE_i/2$。横向性使交换$p,r$所生的$\vct k\cdot\vct\epsilon$项为0。这一化简是能量壳上的跃迁矩元关系（或经幺正长度规范变换），不是对任意态把$p$替成一个经典数。

自旋与轨道磁矩合成$\vct\mu=q(L+g_sS)/(2m)$；对电子是$-\mu_B(L+g_sS)/\hbar$。对称位置张量再分解
\[
 r_ir_j=\frac13r^2\delta_{ij}+\left(r_ir_j-\frac13r^2\delta_{ij}\right),\qquad
 Q_{ij}=q(3r_ir_j-r^2\delta_{ij}).
\]
无源横波$\partial_iE_i=0$消去迹部分，留下$H_{E2}=-Q_{ij}\partial_jE_i/6$；磁部分为$H_{M1}=-\vct\mu\cdot B$。这解释一个长波一阶项为何同时产生反对称的M1和对称无迹的E2，而非两份互不相关的公式。

在体积$V$的一光子态，电场幅度为$\sqrt{\hbar\omega/(2\epsilon_0V)}$。M1的磁场幅度多$1/c$；E2的梯度多$k=\omega/c$，故相应矩元平方为
\[
 |M_{M1}|^2=\frac{\mu_0\hbar\omega}{2V}|\vct\mu_{fi}\cdot(\hat n\times\vct\epsilon)|^2,\qquad
 |M_{E2}|^2=\frac{\hbar\omega^3}{72\epsilon_0Vc^2}|\epsilon_i n_j(Q_{ij})_{fi}|^2.
\]
每偏振每单位能量、立体角的光子态密度是$V\omega^2/[(2\pi)^3\hbar c^3]$。乘黄金规则$2\pi/\hbar$后，角率前因子分别为
\[
 C_{M1}=\frac{\mu_0\omega^3}{8\pi^2\hbar c^3},\qquad
 C_{E2}=\frac{\omega^5}{288\pi^2\epsilon_0\hbar c^5}.
\]
它们还须乘各自的偏振矩元平方，再作偏振和角积分。

横偏振投影$\sum_\lambda\epsilon_i\epsilon_l^*=\delta_{il}-n_in_l$给M1的$|\mu|^2-|\hat n\cdot\mu|^2$，角积分为$8\pi|\mu|^2/3$。E2取对称无迹的复跃迁矩阵$Q$，使用
\[
 \int n_jn_m\dd\Omega=\frac{4\pi}3\delta_{jm},\quad
 \int n_in_ln_jn_m\dd\Omega=\frac{4\pi}{15}
 (\delta_{il}\delta_{jm}+\delta_{ij}\delta_{lm}+\delta_{im}\delta_{lj}),
\]
实际收缩为
\[
 \sum_\lambda\int|\epsilon_i n_jQ_{ij}|^2\dd\Omega
 =\frac{4\pi}3\sum|Q_{ij}|^2
 -\frac{4\pi}{15}\left(2\sum|Q_{ij}|^2+|\Tr Q|^2\right)
 =\frac{4\pi}5\sum|Q_{ij}|^2.
\]
所以总率的系数现在有完整来源：
\[
 \Gamma_{M1}=\frac{\mu_0\omega^3|\vct\mu_{fi}|^2}{3\pi\hbar c^3},\qquad
 \Gamma_{E2}=\frac{\omega^5}{360\pi\epsilon_0\hbar c^5}\sum_{ij}|(Q_{ij})_{fi}|^2.
\]
这分别支撑PS3 Q4的$2S$超精细M1和Q5的$3d\to1s$ E2；该$1/360$采用这里$Q_{ij}=q(3r_ir_j-r^2\delta_{ij})$定义，若改用不含3的无迹张量，系数亦要改。以上长波、电磁弱耦合、原子不反冲近似对应原含时讲义34--37页，自主补出了算符分解与系数积分。M1偶宇称、秩1；E2偶宇称、秩2。相对一个允许E1的典型速率，二者分别受$\alpha^2$和$(ka)^2$抑制，但被禁通道不存在可比较的非零E1速率。

更高阶不宜逐一展开笛卡尔幂。取$\vct X_{jm}=\vct L Y_{jm}/\sqrt{j(j+1)}$，磁型横波为$\vct M_{jm}=j_j(kr)\vct X_{jm}$，电型为$\vct N_{jm}=k^{-1}\nabla\times\vct M_{jm}$；满足横向Helmholtz方程，$j\ge1$。两族对应光子的电/磁多极，场的矢量指标与空间一起旋转，总角动量才是守恒量。它们提供讲义第38页向量球谐讨论的明确构造。

光电离则将束缚末态换成连续电子态。盒归一$V^{-1/2}e^{i\vct p\cdot\vct r/\hbar}$给每单位立体角/能量态密度$Vm p/(2\pi\hbar)^3$，与矩元$1/V$抵消。结果必须含能量阈值和连续波的适用近似；其Hydrogen Fourier积分、截面和角分布见本科“含时微扰增补”。热平衡中吸收$N$、发射$N+1$与原子Boltzmann比共同给$N(\omega)=(e^{\hbar\omega/k_BT}-1)^{-1}$，乘两偏振态密度即Planck谱，不以经典等分配代替。

\originalsection{自然线宽、绝热与几何相位}
原含时讲义40--43页从离散初态耦合连续末态得到自然线宽。为补齐动力学链，取$H=H_0+V$，初态为激发原子加真空$|i\rangle$，末态$f$是低能原子加一光子。先忽略末态之间更高阶耦合，写
\[
 |\psi(t)\rangle=a(t)e^{-iE_it/\hbar}|i\rangle+
 \sum_f b_f(t)e^{-iE_ft/\hbar}|f\rangle,\quad a(0)=1,\quad b_f(0)=0.
\]
Schrödinger方程逐态投影给联立幅度式
\[
 i\hbar\dot a=\sum_fV_{if}e^{i(E_i-E_f)t/\hbar}b_f,\qquad
 i\hbar\dot b_f=V_{fi}e^{i(E_f-E_i)t/\hbar}a.
\]
将第二式从0积分，$b_f(t)=-iV_{fi}\int_0^t a(s)e^{i(E_f-E_i)s/\hbar}\dd s/\hbar$；代回第一式，令$\tau=t-s$，便得到消去末态后的记忆方程
\[
 \dot a(t)=-\int_0^t K(\tau)a(t-\tau)\dd\tau,\qquad
 K(\tau)=\frac1{\hbar^2}\sum_f|V_{fi}|^2e^{i(E_i-E_f)\tau/\hbar}.
\]
这一步尚无指数衰变假设。把末态和写成平滑谱密度$J(E)=\sum_f|V_{fi}|^2\delta(E-E_f)$后，$K$是其Fourier变换；宽而平滑的谱使相关记忆在$\tau_c$内衰减。只有弱耦合$\Gamma\tau_c\ll1$、观察时间$t\gg\tau_c$时，才可把$a(t-\tau)$在记忆窗内换成$a(t)$，并把上限延到$\infty$。这就是Weisskopf--Wigner/Markov近似，有限盒的离散末态一般不满足永久短记忆。

以绝热收敛处方计算时间积分，
\[
 \int_0^\infty e^{i(E_i-E)\tau/\hbar}\dd\tau
 =\pi\hbar\delta(E_i-E)+i\hbar\operatorname{PV}\frac1{E_i-E}.
\]
所以局域幅度方程为
\[
 \dot a=-\left(\frac\Gamma2+\frac{i\Delta E}\hbar\right)a,\quad
 \Gamma=\frac{2\pi}{\hbar}J(E_i),\quad
 \Delta E=\operatorname{PV}\int\frac{J(E)}{E_i-E}\dd E,
\]
从而$a(t)=e^{-\Gamma t/2-i\Delta Et/\hbar}$。总Schrödinger相位对应复能量$E_i+\Delta E-i\hbar\Gamma/2$；$\Gamma$在此是每秒的衰变率，能量宽度为$\hbar\Gamma$。概率$|a|^2=e^{-\Gamma t}$，寿命$1/\Gamma$，不能把幅度半率直接当概率率。

再做一次实际Fourier积分，设$E_c=E_i+\Delta E$，
\[
 \int_0^\infty e^{iEt/\hbar}e^{-iE_ct/\hbar-\Gamma t/2}\dd t
 =\frac1{\Gamma/2-i(E-E_c)/\hbar}.
\]
模平方遂给Lorentz型$[(E-E_c)^2+(\hbar\Gamma/2)^2]^{-1}$。在峰值一半处$|E-E_c|=\hbar\Gamma/2$，故半高全宽$\hbar\Gamma$；若换成光子能量，中心还须减末态原子的能量，真实能移需要初末态都一致处理。把缓变的$J(E)$提出峰宽范围，是这条近似线型成立的另一条件。

该链不涵盖极短时的二次存活或极长时受谱下界控制的尾部。原讲义所提Lamb位移的高频积分还会在简单模型中发散，需更完整QED处理；本册没有用任意截断冒称已算精确Lamb结果。

绝热不是“慢”这一形容词。微分$H(R)|n(R)\rangle=E_n|n\rangle$给异态
$\langle m|\dot n\rangle=\langle m|\dot H|n\rangle/(E_n-E_m)$，控制参数为$\hbar|\langle m|\dot H|n\rangle|/|E_n-E_m|^2\ll1$，且沿路径有隔离能隙。演化为动力相位乘几何相位，后者
\[
 \gamma_n=\int_C\vct{\mathcal A}_n\cdot\dd\vct R,\quad
 \vct{\mathcal A}_n=i\langle n|\nabla_Rn\rangle,\quad
 \mathcal F_{ab}=-2\operatorname{Im}
 \sum_{m\ne n}\frac{\langle n|\partial_aH|m\rangle
 \langle m|\partial_bH|n\rangle}{(E_n-E_m)^2}.
\]
态规范$|n\rangle\to e^{i\chi}|n\rangle$使$\mathcal A\to\mathcal A-\nabla\chi$，曲率不变。局部光滑区域$\mathcal F=\dd\mathcal A$满足Bianchi恒等式，三参数时表现为$\nabla\cdot\vct{\mathcal B}=0$；退化点可如磁单极形成奇点，不能跨过它直接用一个全局势。自旋磁场闭环的$\gamma_m=-m\Omega$见PS4，符号由本章$\mathcal A$定义固定。

\originalsection{群作用、周期势与全同粒子}
对称群在态空间由保持跃迁概率的幺正或反幺正变换作用。连续$U(\epsilon)=e^{-i\epsilon G/\hbar}$若与$H$交换，生成元守恒；同能量空间携带群表示，不可约表示维数提供对称性简并。宇称使$\vct x,\vct p$反号，轴矢量$\vct L$不变；时间反演反幺正，使$i,\vct p,\vct J$反号。Kramers的$\Theta^2=-1$使每个有限维能级偶数维，严格证明见PS6，并非只说“至少两重”。外加固定磁场通常破坏此对称；将磁场本身也反转是比较另一个实验Hamiltonian。

周期势$V(x+a)=V(x)$的平移谱$e^{iKa}$使$\psi(x+a)=e^{iKa}\psi(x)$，即Bloch态$e^{iKx}u_K(x)$。用$\psi,\psi'$组成状态向量，周期单胞转移矩阵$M(E)$的Wronskian守恒给$\det M=1$。允许无限晶格扩展态须两个特征值位于单位圆，因此$|\Tr M|\le2$且$\Tr M=2\cos Ka$；超出即禁带，不以增长解充当可归一束缚态。紧束缚局域轨道$|n\rangle$、最近邻矩元$-t$给$E(K)=E_0-2t\cos Ka$，说明离散能级如何展成能带。零耦合极限是简并孤立阱，不是有限带宽。

全同粒子的空间和自旋合并后才施加交换对称。两个电子的自旋单态搭配对称空间态，三重态搭配反对称空间态。两正交轨道$u,v$的Coulomb直接积分$D$及交换积分$K$使能量$E_{u}+E_v+D\pm K$；$K\ge0$可由Coulomb核的正Fourier变换证明，因此三重态较低。空间分离使$K$趋零，交换效应变小。非相对论理论将Bose/Fermi统计作为输入，不能靠一次360度旋转的图示证明自旋统计定理；后者还依赖相对论局域性等条件。

Hartree乘积变分产生自洽单粒子方程，但电子需Slater行列式与交换项才成为Hartree--Fock。对轨道$\phi_i$，能量$E=\sum_i\langle h_i\rangle+\frac12\sum_{i\ne j}D_{ij}$，其变分势由其他轨道密度构成；把单粒子本征值直接相加会多计相互作用，须减一次每对的直接积分。后续一维两粒子模型展示可解析的自洽解。

\originalsection{Young图与张量空间的完整分解}
先说明这里的表示为何与物理态有关。群表示是$D:G\to GL(\mathcal V)$，满足$D(g_1g_2)=D(g_1)D(g_2)$和$D(e)=I$；量子对称常取保持内积的幺正表示。子空间$W$若对每个$g$有$D(g)W\subseteq W$，称不变；只含$0$和整个空间两个不变子空间时称不可约。有限维幺正表示中不变子空间的正交补也不变，因为$\langle D(g)u,w\rangle=\langle u,D(g)^{-1}w\rangle=0$。于是可以逐块分解，而非只凭矩阵尺寸猜简并。

原对称讲义38--44页用$S_3$和Young表说明置换表示。给三个不同正交标签的张量基$|123\rangle$及其六种排列，$P_{ij}$交换第$i,j$槽。对形状$[2,1]$的表$\left(\begin{smallmatrix}1&2\\3&\end{smallmatrix}\right)$，行对称子为$R=I+P_{12}$、列反对称子为$C=I-P_{13}$。取顺序固定的$Y=CR$，实际作用给
\[
 v=Y|123\rangle=|123\rangle+|213\rangle-|321\rangle-|312\rangle,
 \qquad
 w=P_{23}v=|132\rangle+|231\rangle-|312\rangle-|321\rangle.
\]
这是算符计算，不是把一个图名直接当波函数。逐项交换槽得$P_{12}v=v-w$、$P_{12}w=-w$、$P_{23}v=w$、$P_{23}w=v$，所以$\operatorname{span}(v,w)$在生成整个$S_3$的两换位下不变。$\langle v|v\rangle=\langle w|w\rangle=4$、$\langle v|w\rangle=2$，正交基取$u_1=v/2,u_2=(w-v/2)/\sqrt3$；在此基有
\[
 D(P_{12})=\frac12\begin{pmatrix}1&-\sqrt3\\-\sqrt3&-1\end{pmatrix},\qquad
 D(P_{23})=\frac12\begin{pmatrix}1&\sqrt3\\\sqrt3&-1\end{pmatrix}.
\]
两矩阵平方都是$I$，乘积三次为$I$，满足置换群关系；没有共同不变直线，故这一二维块就是$[2,1]$不可约表示。$Y$本身不是正交投影，需注意$R,C$不交换；本例$Y^2=3Y$，归一或正交化不能任意省略。

相同$U$作用每个张量槽的$U^{\otimes3}$与全部槽置换交换，故它保持这些置换类型。以两维自旋槽为例，将上面形式标签重复代入$1=2=\uparrow,3=\downarrow$，得到两个线性无关、与全对称三重上/下降态正交的$M_S=1/2$混合态；自旋降算符再给$M_S=-1/2$。所以两维置换块伴随一个两维自旋块，
\[
 (\mathbb C^2)^{\otimes3}
 =V^{S_3}_{[3]}\otimes D_{3/2}\ \oplus\ V^{S_3}_{[2,1]}\otimes D_{1/2},\quad
 8=1\cdot4+2\cdot2.
\]
这展示Schur--Weyl分解中两个互相交换的群分别作用哪一因子；不把$S_3$维数与自旋维数混为同一个重数。下面给一般计数公式及后续作业的具体验算，一般hook长度/Schur--Weyl定理的完整证明不在此例中声称完成。

置换群$S_N$不可约表示由分拆$\lambda\vdash N$标记。标准Young表在行和列都严格递增，其数目为$f^\lambda=N!/\prod_{(i,j)\in\lambda}h_{ij}$，$h_{ij}$是该格右方与下方格数加1；例如$S_3$的$[3],[2,1],[1^3]$维数$1,2,1$，平方和$6=3!$。行对称子与列反对称子构造对应子空间，操作顺序和归一需保持一致；随意将所有表的向量列在一起可能过完备。

单粒子内部空间为$\mathbb C^k$时，两种互相交换的作用产生Schur--Weyl分解
\[
 (\mathbb C^k)^{\otimes N}
 =\bigoplus_{\substack{\lambda\vdash N\\\ell(\lambda)\le k}}
 V^{S_N}_\lambda\otimes V^{U(k)}_\lambda,\qquad
 \dim V^{U(k)}_\lambda=\prod_{(i,j)\in\lambda}\frac{k+j-i}{h_{ij}}.
\]
因为一列超过$k$个向量的反对称积为零，只有至多$k$行。$SU(k)$中完整$k$格列是行列式单态，可删去；要区分$S_N$表示维数$f^\lambda$与$SU(k)$表示维数，不可把两者都称“重数”。PS8逐项给$(\mathbb C^2)^5$和$(\mathbb C^3)^4$的维数验算。

张量积按Littlewood--Richardson规则：将第二图逐行标字，添加到第一图，所得格形仍为Young图；每行弱递增、每列严格递增，右至左、上至下读字的任一前缀中$a$不少于$b$、$b$不少于$c$。不同合法填法给重数。于是$SU(3)$有$3\otimes3=6\oplus\bar3$、$3^{\otimes3}=10\oplus8\oplus8\oplus1$，维数$27$。这不是任意$S_N$两个不可约表示张量积的同一规则，后者是Kronecker系数问题。

原讲义以$p^3$原子与三夸克为应用。电子空间/自旋表示须组合出完全反对称；基态重子颜色为反对称单态、空间为对称，故自旋--味部分需对称。PS8给完整$p^3$的$^2P$基和质子味--自旋波函数，并给$d^3$全部项符号；避免只列Young图而不说明物理态。

\originalsection{散射边界、幺正与极点}
相对坐标将两体弹性散射化成约化质量$\mu$的势散射。定义$\psi\sim e^{ikz}+f(\theta,\phi)e^{ikr}/r$，流比给$\dd\sigma/\dd\Omega=|f|^2$。只对等入/出速度的弹性散射成立；非弹性须另乘速度比。Lippmann--Schwinger式
$|\psi^{(+)}\rangle=|\vct k\rangle+(E-H_0+i0)^{-1}V|\psi^{(+)}\rangle$
由$+i0$固定出射边界；位置Green函数$-e^{ikr}/(4\pi r)$满足$(\nabla^2+k^2)G=\delta$。取一次迭代得
\[
 f_B(\vct k',\vct k)=-\frac{\mu}{2\pi\hbar^2}
 \int e^{-i(\vct k'-\vct k)\cdot\vct r}V(r)\dd^3r.
\]
这是弱势/高能的第一Born近似。实势下它在前向常为实，不能据此说精确光学定理失效：总截面二阶而前向虚部也要二阶计算。精确概率流的大球面积分中，入射--散射的前向干涉抵消散射总流，得到$\sigma_{\rm tot}=4\pi\operatorname{Im}f(0)/k$。

中心势分波$S_l=e^{2i\delta_l}$；径向实解在势外为$\cos\delta_l\,j_l(kr)-\sin\delta_l\,n_l(kr)$。设边界$r=R$的对数导数$\beta_l=Ru_l'/u_l-1$，匹配给
\[
 \tan\delta_l=\frac{kRj_l'(kR)-\beta_lj_l(kR)}
 {kRn_l'(kR)-\beta_ln_l(kR)},\quad
 f=\frac1k\sum_l(2l+1)e^{i\delta_l}\sin\delta_lP_l(\cos\theta),\quad
 \sigma=\frac{4\pi}{k^2}\sum_l(2l+1)\sin^2\delta_l.
\]
每分波上限$4\pi(2l+1)/k^2$来自幺正。$kf_l=e^{i\delta_l}\sin\delta_l$在复平面位于中心$i/2$、半径$1/2$的圆上；吸收时$|S_l|<1$位于圆内，差额是反应截面。

短程势低能高分波受离心势垒抑制，$\delta_l=O(k^{2l+1})$。定义零能$s$波$u\propto r-a_s$，所以$\delta_0\simeq-ka_s$，$\sigma\to4\pi a_s^2$。更精确$k\cot\delta_0=-1/a_s+r_ek^2/2+\cdots$，$f_0=[k\cot\delta_0-ik]^{-1}$。若$a_s\gg R>0$，极点$k\simeq i/a_s$给浅束缚能$-\hbar^2/(2\mu a_s^2)$。负散射长度可对应虚态，不能据散射长度正负一概断定势的吸引/排斥。硬球$a_s=R$，高能总截面趋$2\pi R^2$：几何反射与前向衍射各贡献$\pi R^2$。六个指定波数的图见PS11。

共振的$S_l$可写成背景相位乘$(E-E_R-i\Gamma/2)/(E-E_R+i\Gamma/2)$，在下半能量平面有极点；忽略背景且$k$近似常时是Breit--Wigner峰$\sigma_l\simeq4\pi(2l+1)k^{-2}(\Gamma/2)^2/[(E-E_R)^2+(\Gamma/2)^2]$。波包群延迟$2\hbar\,\dd\delta_l/\dd E$体现暂困时间，但不等于每个能量都统一的寿命。讲义最后还将Coulomb问题化为合流超几何方程；Coulomb长程渐近含$-\eta\ln(2kr)$，不能套用上述短程自由渐近。这里补出计算桥接。取$\rho=kr$、$\eta=\mu Z_1Z_2g/(\hbar^2k)$、$g=e^2/(4\pi\epsilon_0)$，径向式$u''_\rho+[1-2\eta/\rho-l(l+1)/\rho^2]u=0$。设$u=\rho^{l+1}e^{i\rho}v(-2i\rho)$，则$xv''+(2l+2-x)v'-(l+1+i\eta)v=0$。原点正则支是Kummer $M(l+1+i\eta,2l+2,x)$；其大$|x|$的两支渐近给$u\sim\sin[\rho-\eta\ln(2\rho)-l\pi/2+\sigma_l]$，$\sigma_l=\arg\Gamma(l+1+i\eta)$。Kummer渐近为此处明示的特殊函数输入，不能删去对数相位后称自由球波。

Coulomb径向出射波应带$e^{i[kr-\eta\ln(2kr)]}$；入射也含长程畸变相位，不是前文短程势的纯平面波渐近。将共同的径向对数相位提出后，Coulomb分波相移给
\[
 S_l=e^{2i\sigma_l}=\frac{\Gamma(l+1+i\eta)}{\Gamma(l+1-i\eta)}.
\]
为实际完成振幅求和，令$x=\cos\theta$，先插入Abel因子$0<r<1$：
\[
 f_r(x)=\frac1{2ik}\sum_{l=0}^\infty(2l+1)(S_l-1)P_l(x)r^l.
\]
下面只取非前向$x<1$；恒等部分由Legendre生成函数给
$\sum(2l+1)P_l(x)r^l=(1-r^2)/(1-2rx+r^2)^{3/2}$，在固定$x<1$时$r\to1$趋0。它在前向是集中的分布，不能在整个角域当普通函数删除。

计算剩余$\Sigma(x)=\sum(2l+1)S_lP_l(x)$。暂令$a=i\eta$满足$-1<\operatorname{Re}a<0$，使用Euler Beta积分
\[
 S_l=\frac1{\Gamma(-2a)}\int_0^1t^{l+a}(1-t)^{-1-2a}\dd t.
\]
先保留$r<1$交换积分与和，再在安全收敛带例如$-1<\operatorname{Re}a<-3/4$及$-1<x<1$取$r\to1$；随后作解析延拓至实际实$\eta$的Abel边界值。将同一个生成函数代入，得到真正可算的一维积分
\[
 \Sigma(x)=\frac1{\Gamma(-2a)}\int_0^1
 \frac{t^a(1-t)^{-1-2a}(1-t^2)}{[1-2xt+t^2]^{3/2}}\dd t.
\]
关键代换是$y=(1-t)^2/t$、$C=2(1-x)>0$。因为
$1-2xt+t^2=t(y+C)$、$\dd y=-(1-t^2)t^{-2}\dd t$，所有$t$的幂恰好抵消，余下
\[
 \Sigma(x)=\frac1{\Gamma(-2a)}\int_0^\infty
 y^{-a-1/2}(y+C)^{-3/2}\dd y
 =C^{-a-1}\frac{\Gamma(1/2-a)\Gamma(1+a)}{\Gamma(3/2)\Gamma(-2a)}.
\]
最后以Gamma倍角式
$\Gamma(-2a)=2^{-2a-1}\Gamma(-a)\Gamma(1/2-a)/\sqrt\pi$
和$\Gamma(1-a)=-a\Gamma(-a)$化简，记$h=\sin^2(\theta/2)$、$C=4h$，则
\[
 \Sigma(x)=-a\frac{\Gamma(1+a)}{\Gamma(1-a)}h^{-a-1}.
\]
因此非前向极限从上述积分、而不是一句“正则求和”给出
\[
 f_C(\theta)=-\frac{\eta}{2k\sin^2(\theta/2)}
 e^{-2i\eta\ln\sin(\theta/2)+2i\sigma_0},\qquad
 \frac{\dd\sigma}{\dd\Omega}=\frac{\eta^2}{4k^2\sin^4(\theta/2)}.
\]
这里$\theta=\pi$由非前向连续极限定义，$\eta=0$由自由极限定义；$\theta=0$仍发散。实际实$\eta$的级数未被当成绝对收敛普通和；所用的是Abel边界值及带Coulomb对数相位的长程散射态。这一推导将8.322散射讲义末页的Coulomb微分方程和相移补成可核查的振幅计算桥接，新增积分细节为自主展开。幅度模给Rutherford，精确相位却不是第一Born实振幅。

吸引时解析延拓$k=i\kappa$，不增长的多项式条件$l+1+i\eta=-n_r$给$\kappa=\mu|Z_1Z_2|g/[\hbar^2(n_r+l+1)]$，从而回Hydrogen束缚谱；排斥没有该正$\kappa$支。一般短程光学定理的总截面积分在Coulomb前向发散，需屏蔽/有限角分辨。核/原子形状因子在PS9、11中另作。

\originalsection{课历与原页覆盖定位}
课历从含时微扰和辐射，经过绝热、对称、全同粒子/原子/群表示，再到散射、相对论量子力学；官方网页课历按讲次分组：1--8含时/辐射，9--16对称与全同粒子，17--22散射，23--26相对论；未给逐讲具体日期。另取得的3页PDF含当年授课/习题课时间及前学期回顾，不能冒称逐日课历。下章接续Dirac理论。为避免把合本页数当内容完成率，本次只按实际主题和缺件记录覆盖：
\begin{longtable}{p{.26\textwidth}p{.17\textwidth}p{.47\textwidth}}
\toprule 合本 & 原页范围 & 中文定位\\\midrule\endhead
含时/辐射 & 1--17 & 本节黄金规则、上一分编旋转系/Dyson、本科含时辐射；氨分子补充。\\
 & 18--27 & 最小耦合、偶极、光电离态密度；20--23为重复原页。\\
 & 28--39 & 场量子化、E1/M1/E2、向量球谐、热平衡，PS3。\\
 & 40--51 & 自然线宽、绝热条件、Berry曲率，PS4。\\
对称 & 1--16 & 群表示、宇称、反幺正/Kramers；上一分编与PS6。\\
 & 17--26 & Bloch、紧束缚、转移矩阵和能带，PS6。\\
 & 27--37 & 全同粒子、交换、Hartree；PS7。\\
 & 38--53 & Young表、Schur--Weyl、张量积、原子与重子；PS8。\\
散射 & 1--10 & 弹性定义、出射Green、Born、光学定理；PS9。\\
 & 11--23 & 球Bessel、分波归一/匹配、幺正圆、硬球与衍射；PS10--11。\\
 & 24--30 & 散射长度、束缚极点、共振、Coulomb边界；本章及本科散射。\\
\bottomrule\end{longtable}
